\documentclass[11pt]{article}
\usepackage[final]{acl}
\usepackage{times}
\usepackage{latexsym}
\usepackage{booktabs}
\usepackage{multirow}
\usepackage[T1]{fontenc}
\usepackage[utf8]{inputenc}
\usepackage{microtype}
\usepackage{inconsolata}
\usepackage{xeCJK}
\setCJKmainfont{Songti SC}

\title{Four Repositories, One List: The Open HSK 3.0 Datasets Encode\\
       a Document the Examination No Longer Uses}

\author{Alvaro Serrano \\
  \textit{Independent researcher} \\
  \texttt{alvaro.serrano.gp@gmail.com} \\
  \texttt{orcid.org/0009-0006-4701-9026}}

\begin{document}
\maketitle

\begin{abstract}
Tools that show a learner an HSK level almost never derive it from the issuing
body's PDF. They take it from a small number of open datasets. We audit the five
most-starred, comparing each word by word against the three documents in
circulation: the 2012 syllabus, the GF0025-2021 national standard, and the 2025
examination syllabus in force since July 2026. \textbf{Four of the five encode
GF0025-2021, and only one encodes the examination.} Worse for anyone treating
their agreement as corroboration, the four are not independent: across every
pair, over roughly 10{,}940 shared words, they disagree about the level of
\textbf{not one word}. They are one list. The derivation is documented rather
than hidden --- one repository credits another as its source, and that one links
the March 2021 ministry PDF it was OCR'd from --- so the ecosystem is not
careless; it is \emph{singly rooted}, which is a different and more durable
problem. A consumer three hops downstream sees only ``HSK 3.0.'' The audit also yields
an external replication we did not set out to obtain. Comparing a community 2021
list to a community 2025 list --- neither ours, from different maintainers ---
reproduces the disagreement rate reported in \citet{serrano2026whichhsk} to
within \textbf{0.01 percentage points} (41.49\% against 41.48\%), on data with no
dependency on our extraction. We release the audit, which re-runs against the
live repositories in one command.
\end{abstract}

\section{Nobody reads the PDF}

The 2025 examination syllabus is distributed as a 406-page PDF. The 2021
national standard is a book. Extracting either into a machine-readable list is a
week of careful work, and almost nobody does it: a reader, an app, or a
levelling script takes the list from GitHub, where a handful of open datasets
have accumulated thousands of stars between them and are reused far beyond any
audit trail.

That makes those datasets the field's actual standard, whatever the issuing
bodies publish. If they encode one document while learners are examined against
another, the error is not any one vendor's. It is inherited.

\citet{serrano2026whichhsk} established that ``HSK 3.0'' names two official
documents which assign different levels to 41.5\% of the vocabulary they share.
This paper asks the obvious follow-up, which is not about the documents at all:
\emph{which one is actually in the software?}

\section{Method}

We take the five most-starred open HSK 3.0 word lists (Table~\ref{tab:audit}) and
compare each, word by word, against all three documents. Agreement is computed
over the words a list and a document \emph{share}, so an incomplete list is not
penalised for being incomplete; what is measured is whether it assigns the same
level. Levels 7, 8 and 9 are one undifferentiated band in the source documents
and are treated as one here.

Two normalisations are applied identically to every source, because the official
tables carry both and a list that reproduces them faithfully still has to be
matched: part-of-speech annotations in full-width parentheses (白（形）) are
stripped, and alternates written with a full-width bar (爸爸｜爸) are reduced to
the headword.

\paragraph{The control.} One repository publishes both documents, in folders
named \texttt{New HSK (2021)} and \texttt{New HSK (2025)}. An audit that could
not tell those apart would be worthless. Each matches its own label at 100.0\%,
and each matches the other document at 58.5\%. The method separates them.

\paragraph{What we do not redistribute.} Every list is fetched from its own
repository under its own licence, used to compute summary statistics, and not
republished. The audit script ships; the data does not.

\section{Results}

\begin{table*}[t]
\centering\small
\begin{tabular}{@{}lrrrrr@{}}
\toprule
\textbf{List} & \textbf{Stars} & \textbf{Words} & \textbf{HSK 2.0} &
\textbf{GF0025-2021} & \textbf{2025 syllabus} \\
\midrule
\texttt{krmanik/HSK-3.0 (2025)}          & 356 & 10{,}900 & 48.9\% & 58.5\% & \textbf{100.0\%} \\
\texttt{krmanik/HSK-3.0 (2021)}          & 356 & 10{,}943 & 18.2\% & \textbf{100.0\%} & 58.5\% \\
\texttt{drkameleon/complete-hsk-vocabulary} & 294 & 10{,}969 & 18.2\% & \textbf{100.0\%} & 58.5\% \\
\texttt{elkmovie/hsk30}                  & 115 & 10{,}946 & 18.2\% & \textbf{100.0\%} & 58.5\% \\
\texttt{ivankra/hsk30}                   &  80 & 10{,}946 & 18.2\% & \textbf{100.0\%} & 58.5\% \\
\bottomrule
\end{tabular}
\caption{Agreement with each document, over the words the list and the document
share (between 9{,}694 and 10{,}942 in every cell above). Star counts read
5 September 2026.}
\label{tab:audit}
\end{table*}

\subsection{One of five encodes the examination}

Table~\ref{tab:audit} is the result. Four of the five lists match GF0025-2021
exactly. The single list that matches the 2025 examination syllabus is the one
whose own directory name says 2025.

GF0025-2021 is the national \emph{grading standard}. It is not wrong, and a
programme teaching to the national standard is right to use it. But learners
sitting the HSK from July 2026 are examined against the other document, and a
tool that labels a word ``HSK 4'' from one of these lists is not making a
statement about the examination its user is preparing for.

\subsection{The four are the same list}

\begin{table}[t]
\centering\small
\begin{tabular}{@{}lrr@{}}
\toprule
\textbf{Pair} & \textbf{Shared} & \textbf{Differ} \\
\midrule
krmanik(2021) vs drkameleon & 10{,}935 & \textbf{0} \\
krmanik(2021) vs elkmovie   & 10{,}935 & \textbf{0} \\
krmanik(2021) vs ivankra    & 10{,}935 & \textbf{0} \\
drkameleon vs elkmovie      & 10{,}938 & \textbf{0} \\
drkameleon vs ivankra       & 10{,}938 & \textbf{0} \\
elkmovie vs ivankra         & 10{,}942 & \textbf{0} \\
\addlinespace
any of the four vs krmanik(2025) & $\sim$9{,}700 & $\sim$4{,}025 \\
\bottomrule
\end{tabular}
\caption{Pairwise level disagreements among the audited lists.}
\label{tab:pairs}
\end{table}

Table~\ref{tab:pairs} is the finding that matters more. Across all six pairs
among the four GF0025-2021 lists, over roughly 10{,}940 shared words each, the
number of words assigned a different level is \textbf{zero}.

Not ``low''. Zero. The only differences anywhere were four entries where one list
writes a homograph index as a superscript (称¹) and another as a digit (称1).

\textbf{This is not a discovery of hidden copying.} The chain is documented.
\texttt{drkameleon} lists its sources in its README and credits
\texttt{elkmovie/hsk30} for the HSK 3.0 wordlist; \texttt{elkmovie} states in its
own header that it was OCR'd with Pleco OCR and links the March 2021 Ministry of
Education announcement it was taken from. \texttt{ivankra} documents its row
count as 11{,}092 --- the GF0025-2021 total --- and its ID scheme as the index
from that PDF. Every maintainer said what they did.

The finding is structural rather than editorial. The open Chinese-proficiency
ecosystem is \emph{singly rooted}: essentially all of it descends, by a
disclosed path, from one OCR pass over one PDF published in March 2021. That is
efficient, and it is fragile in two specific ways.

First, \textbf{corroboration is unavailable}. A practitioner who checks a word
against three HSK lists and finds them in agreement has learned nothing about the
standard; they have learned that one file propagated. Agreement across these
lists cannot detect an OCR error, and the root list says it was ``not
extensively proofread.''

Second, \textbf{the root is now the wrong document}. Nothing in the chain is a
mistake --- three of the four lists predate the 2025 syllabus entirely --- but
the ecosystem inherits its currency from its root, and the root is the 2021
standard. Updating one repository does not update the field; updating the field
means each downstream consumer independently noticing that a second document now
exists.

\subsection{Nobody says which}

Attribution within the chain is good. What degrades is what a \emph{consumer}
of these lists can see.

At the root, \texttt{elkmovie} links the exact ministry announcement its list
came from, dated March 2021. That is unambiguous provenance. One hop downstream,
\texttt{drkameleon} credits \texttt{elkmovie} but presents its own levels under
the heading ``HSK 3.0'' --- a name that, since November 2025, denotes two
documents. Two hops downstream, an application that imports
\texttt{complete.min.json} and renders \texttt{n4} as ``HSK 4'' has no surviving
reference to March 2021 at all, and neither does its user.

So the information is not lost through negligence; it is lost through
\emph{distance}. Provenance that lives in a README does not travel with the
data. Only the repository that explicitly split its output into \texttt{New HSK
(2021)} and \texttt{New HSK (2025)} folders carries the distinction in a form
that survives being copied, and that is the one list of the five that a tool can
use without inheriting the ambiguity.

The remedy is correspondingly small: a version field \emph{inside} the data,
naming the source document and its date, so that the answer travels with the
answer.

\section{An external replication}
\label{sec:replication}

This was not the point of the audit, and it is the strongest thing in it.

\begin{table}[t]
\centering\small
\begin{tabular}{@{}lrrr@{}}
\toprule
\textbf{Lists compared} & \textbf{Shared} & \textbf{Differ} & \textbf{Rate} \\
\midrule
Community datasets & 9{,}700 & 4{,}025 & \textbf{41.49\%} \\
Our own extraction & 9{,}698 & 4{,}023 & \textbf{41.48\%} \\
\bottomrule
\end{tabular}
\caption{The headline finding of \citet{serrano2026whichhsk}, recomputed from
data we did not produce. \texttt{drkameleon} supplies the 2021 side and
\texttt{krmanik} the 2025 side; they are different maintainers.}
\label{tab:replication}
\end{table}

Take a community list that encodes GF0025-2021 and a community list that encodes
the 2025 syllabus, from different maintainers, and compare those two to each
other. Neither has any dependency on our extraction of either document.

The rates differ by \textbf{0.01 percentage points} (Table~\ref{tab:replication}).

This is stronger evidence than the original measurement. An extraction error on
our side could manufacture a 41.5\% disagreement; it could not manufacture the
same figure a second time in data extracted by other people from the same two
PDFs. The finding is a property of the documents, not of our pipeline.

It also cuts the other way, and we state it as such: the community lists
reproduce our number because they and we are reading the same two documents
correctly. Table~\ref{tab:replication} validates our extraction as much as it
validates the claim.

\section{Limitations}

We audit five lists. There are more, and the selection is by GitHub stars, which
measures visibility rather than deployment --- we cannot observe which list any
given product actually ships, and we make no claim about any named product.

Star counts are a snapshot and will drift; the audit re-runs and reports the
counts it read.

Agreement over shared words does not test coverage. A list could match a
document perfectly on every word it contains and still omit a third of it. We
report list sizes so this is visible, but sizes are not a completeness audit.

The zero-disagreement result establishes that the four lists are the same
artifact, and the repositories' own READMEs establish the direction of
derivation for the pair that documents it. We do not assert a derivation path
where none is stated.

Finally, this measures word lists. A tool's text-level output also depends on
segmentation, a coverage threshold and proper-noun handling, none of which is
audited here. The narrower claim --- about the list, which a published document
is supposed to fix --- is the one we make.

\section*{Availability}

\texttt{scripts/audit\_wordlists.py} in the \texttt{hsk30} repository reproduces
every figure here, fetching each list live from its own repository. No audited
list is redistributed. A companion,
\texttt{scripts/standard\_fingerprint.py}, emits a twelve-word probe that
identifies which of the three documents any black-box tool implements, for the
common case where the list is not published at all.

\section*{Ethics and data statement}

No personal data and no human subjects. All audited datasets are public and
openly licensed; each is used to compute summary statistics and none is
republished. We name repositories because the finding is about what a user of a
public artifact can know from it, and naming is necessary for the claim to be
checkable --- but no finding here implies error or bad faith by any maintainer.
Three of the four lists predate the document they are measured against, every
maintainer documented what they did, and an earlier draft of this paper
mischaracterised two of them by not reading their READMEs before writing about
their data.
Generative AI was used in writing the audit code and in drafting this paper; the
framing, the design decisions and every claim are the author's, who reviewed all
output. That review is what caught two parser defects --- an OCR'd section header
and a level written \texttt{7-9} --- each of which had silently produced a wrong
result about somebody else's repository before it was corrected.

\bibliography{refs}

\end{document}
